\chapter{Numerical Linear Algebra}

\section{Linear Systems and Direct Methods}

A system of $m$ linear equations in $n$ unknowns can be written as
\[
  A\vect{x}=\vect{b},
\]
where $A=[a_{ij}]\in\R^{m\times n}$ is the coefficient matrix,
$\vect{x}=[x_1,\ldots,x_n]^T$ is the vector of unknowns, and
$\vect{b}=[b_1,\ldots,b_m]^T$ is the right-hand side.  In this course we are mainly interested in square systems, $m=n$.

We have previously learned a few direct methods to simplify linear systems so that an exact solution can easily be found. Gaussian elimination and matrix factorisation are such methods.

\subsection{Gaussian Elimination}

Gaussian elimination transforms the augmented matrix $[A\mid \vect b]$ to row-echelon form (i.e. into an upper-triangular system), or in matrix notation $[U\mid \vect b^*]$. For a pivot in row $k$, the entries below the pivot are removed using elementary row operations as follows:
\[
 R_j-m_{jk}R_k \rightarrow R_j,
 \qquad
 m_{jk}=\frac{a_{jk}^{(k)}}{a_{kk}^{(k)}},
 \qquad j=k+1,\ldots,n.
\]
For reference, $a_{ij}^{(k)}$ refers to the element in the $i$-th row and $j$-th column while eliminating entries below the pivot in the $k$-th row.
Once an upper-triangular system $U\vect{x}=\vect{b^*}$ is obtained, backward substitution gives
\[
 x_i=\frac{1}{u_{ii}}
 \left(c_i-\sum_{j=i+1}^{n}u_{ij}x_j\right),
 \qquad i=n,n-1,\ldots,1.
\]

\begin{algorithm}[Gaussian elimination]
\begin{enumerate}
  \item Form the augmented matrix $[A\mid \vect b]$.
  \item For $k=1,\ldots,n-1$, use row operations to eliminate every entry below the pivot $a_{kk}^{(k)}$.
  \item Solve the resulting upper-triangular system by backward substitution.
\end{enumerate}
\end{algorithm}

\begin{example}
Solve the following system of equations with Gaussian elimination:
\[
\begin{array}{ccccccccc}
x_1&+&x_2&+& & &3x_4&=&4,\\
2x_1&+&x_2&-&x_3&+&x_4&=&1,\\
3x_1&-&x_2&-&x_3&+&2x_4&=&-3,\\
-x_1&+&2x_2&+&3x_3&-&x_4&=&4.
\end{array}
\]
\end{example}

\begin{solution}[8.0cm]
Form the augmented matrix:
\[
\left[
\begin{array}{rrrr|r}
1&1&0&3&4\\
2&1&-1&1&1\\
3&-1&-1&2&-3\\
-1&2&3&-1&4
\end{array}
\right].
\]
Then,
\[
\left[
\begin{array}{rrrr|r}
1&1&0&3&4\\
2&1&-1&1&1\\
3&-1&-1&2&-3\\
-1&2&3&-1&4
\end{array}
\right]
\xrightarrow[
  {\displaystyle
    \begin{array}{l}
      R_4+R_1\to R_4
    \end{array}}
]{
  \begin{array}{l}
    R_2-2R_1\to R_2\\
    R_3-3R_1\to R_3
  \end{array}
}
\left[
\begin{array}{rrrr|r}
1&1&0&3&4\\
0&-1&-1&-5&-7\\
0&-4&-1&-7&-15\\
0&3&3&2&8
\end{array}
\right]
\]
\[
    \xrightarrow[
  {\displaystyle
    \begin{array}{l}
      R_4+3R_2\to R_4
    \end{array}}
]{{\displaystyle
    \begin{array}{l}
      R_3-4R_2\to R_3
    \end{array}}
}
\left[
\begin{array}{rrrr|r}
1&1&0&3&4\\
0&-1&-1&-5&-7\\
0&0&3&13&13\\
0&0&0&-13&-13
\end{array}
\right].
\]
Backward substitution yields
\[
   x_4=\frac{-13}{-13}=1,
 \quad x_3=\frac{13-13(1)}{3}=0,
\]\[
 \quad x_2=\frac{-7+5(1)+1(0)}{-1}=2,
 \quad x_1=\frac{4-3(1)-0(0)-2}{1}=-1.
\]
Thus
\[
 \vect{x}=[-1,2,0,1]^T.
\]
\end{solution}

In most cases, Gaussian elimination is the preferred way to solve a linear system $A\vect{x}=\vect{b}$. But there are two problems with Gaussian Elimination, which we will discuss in the following subsections.

\subsection{Pivoting Strategies}

We have previously abstracted the steps for Gaussian Elimination. However, if we were to follow the steps blindly, we may run into a few problems (division by zero or numerical instabilities).

A zero pivot prevents the next elimination step (since it would cause division by zero if we were to eliminate blindly).  Even a nonzero but very small pivot can magnify round-off error because the multipliers $m_{jk}$ become large.

\begin{example}[Zero pivot]
Solve
\[
\begin{cases}
 0x_1+x_2=1 \\
 x_1+x_2=2
\end{cases}
\]
\end{example}

\begin{solution}[2.4cm]
After interchanging the two equations, the system becomes
\[
 x_1+x_2=2,
 \qquad x_2=1,
\]
and thus $x_2=1$ and $x_1=1$.
\end{solution}

\begin{example}[Small pivot]
Consider
\[
\begin{cases}
 \epsilon x_1+x_2=1\\
 x_1+x_2=2
\end{cases}
\]
where $\epsilon$ is very small.  We will see why elimination without pivoting may be unreliable.
\end{example}

\begin{solution}[3.0cm]
Using the first equation as pivot, we will perform the row operation $\displaystyle R_2-\frac{1}{\epsilon}R_1\to R_2$. Notice that the multiplier $1/\epsilon$ is very large.  The second equation becomes
\[
 \left(1-\frac1\epsilon\right)x_2
 =2-\frac1\epsilon
 \implies
 x_2 = \frac{2-\epsilon^{-1}}{1-\epsilon^{-1}}\approx 1.
\]
Substituting back into the first equation, we will get:
\[
    x_1 = \epsilon^{-1}(1-x_2) \approx \epsilon^{-1}(1-1) = 0.
\]
It is clear that this solution became unreliable when we executed Gaussian Elimination numerically. 

\noindent
(Substituting our solution $x_1 = 0, x_2 = 1$ into Equation 2 yields a contradiction.)
\end{solution}

\begin{definition}[Partial pivoting]
At a given step $k$, partial pivoting exchanges row $k$ (the original pivot row) with a row below it with the largest leading element. Mathematically, we want a row $p\ge k$ for which
\[
 |a_{pk}^{(k)}|=
 \max_{k\le i\le n}|a_{ik}^{(k)}|.
\]
Partial pivoting is applied in every step of elimination so that we can get a safe and accurate solution.
\end{definition}

\newpage
\begin{example}
Apply Gaussian elimination with partial pivoting to solve
\[
\begin{cases}
\epsilon x_1 + x_2 = 1\\
x_1 + x_2 = 2
\end{cases}
\]
when \(\epsilon\) is small.
\end{example}

\begin{solution}[3.5cm]
We see in the first column that $|1| > |\epsilon|$ which is under the pivot element. Thus exchange the two equations:
\[
\begin{cases}
 x_1+x_2=2\\
 \epsilon x_1+x_2=1
\end{cases}
\]
Then applying row operation $R_2 - \epsilon R_1 \to R_2$ gives
\[
 (1-\epsilon)x_2=1-2\epsilon,
\]
and hence
\[
 x_2=\frac{1-2\epsilon}{1-\epsilon}\approx1,
 \qquad
 x_1=2-x_2\approx1.
\]
\end{solution}

\subsection{Matrix Inversion}
If $A$ is nonsingular, its inverse can be computed by applying Gauss-Jordan elimination (i.e. reducing a matrix to r.r.e.f., reduced row-echelon form) to the block matrix $[A\mid I]$ until the left block becomes $I$:
\[
 [A\mid I]\longrightarrow[I\mid A^{-1}].
\]
However, when considering a square matrix of order $n$, finding the inverse using this method would mean applying Gauss-Jordan Elimination to an $n \times 2n$ augmented matrix.

\medskip
The issue with speed could be resolved for certain classes of problems. The table below shows the operation counts for Gaussian Elimination:
\[
\begin{array}{l|cc}
    \hline
    & \text{Multiplication/Divisions} & \text{Additions/Subtractions}\\
    \hline
    \text{Elimination Steps}& \rule{0pt}{4.0ex}\dfrac{2n^3+3n^2-5n}{6} &\rule{0pt}{4.0ex}\dfrac{n^3-n}{3} \\[1.5ex]
    \hline
    \text{Backward Substitution}& \rule{0pt}{4.0ex}\dfrac{n^2+n}{2}&\rule{0pt}{4.0ex}\dfrac{n^2-n}{2}\\[1.5ex]
    \hline
\end{array}
\]

From here, we see that the number of operations to do backward substitution is significantly smaller than doing the elimination steps.
Just considering a linear system with a $1000 \times 1000$ matrix, we would need 667,165,500 operations for elimination, but only 1,000,000 operations ($0.15\%$) to perform backward substitution.

Hence, for efficiency, when we compute the inverse for a matrix $A$, we will reduce $[A|I]$ to $[U|Y]$ where $U$ is an upper triangular matrix (or as we know the r.e.f. of $A$) and $Y$ is the matrix we get after performing the same elimination steps on $I$.
Then, we will individually obtain the columns of $A^{-1}$ through solving $U\vect{x}_i=\vect{y}_i$ where $\vect{y}_i$ denotes the $i$-th column of $Y$.

\begin{example}
Use Gaussian Elimination followed by backward substitution to find the inverse of
\[
 A=\begin{bmatrix}
 1&2&-1\\
 2&1&0\\
 -1&1&2
 \end{bmatrix}.
\]
\end{example}

\begin{solution}[5.0cm]
\[
[A|I]=
\left[\begin{array}{ccc|ccc}
1 & 2 & -1 & 1 & 0 & 0\\
2 & 1 & 0 & 0 & 1 & 0\\
-1 & 1 & 2 & 0 & 0 & 1\\
\end{array}\right]
\xrightarrow[
{\displaystyle
\begin{array}{l}
R_3+R_1\to R_3
\end{array}}]
{{\displaystyle
\begin{array}{l}
R_2-2R_1\to R_2
\end{array}}
}
\left[\begin{array}{ccc|ccc}
1 & 2 & -1 & 1 & 0 & 0\\
0 & -3 & 2 & -2 & 1 & 0\\
0 & 3 & 1 & 1 & 0 & 1
\end{array}\right]
\]
\[
\xrightarrow{\begin{array}{l}\displaystyle{R_3+R_2\to R_3}\end{array}}
\left[\begin{array}{ccc|ccc}
1 & 2 & -1 & 1 & 0 & 0 \\
0 & -3 & 2 & -2 & 1 & 0\\
0 & 0 & 3 & -1 & 1 & 1
\end{array}\right]=[U|Y]
\]
\medskip
\[
\begin{bmatrix}
1 & 2 & -1\\
0 & -3 & 2\\
0 & 0 & 3
\end{bmatrix}
\begin{bmatrix}
x_1 \\ x_2 \\ x_3
\end{bmatrix}
=
\begin{bmatrix}
1 \\ -2 \\ -1
\end{bmatrix}
\]
\[
\implies
x_3 = -\frac13, \; x_2 = \frac{-2-2(-\frac13)}{-3} = \frac49, \; x_1 = \frac{1+1(-\frac13)-2(\frac49)}{1}=-\frac29.
\]

\medskip

\[
\begin{bmatrix}
1 & 2 & -1\\
0 & -3 & 2\\
0 & 0 & 3
\end{bmatrix}
\begin{bmatrix}
x_1 \\ x_2 \\ x_3
\end{bmatrix}
=
\begin{bmatrix}
0 \\ 1 \\ 1
\end{bmatrix}
\]
\[
\implies
x_3 = \frac13, \; x_2 = \frac{1-2(\frac13)}{-3} = -\frac19, \; x_1 = \frac{0+1(\frac13)-2(-\frac19)}{1}=\frac59.
\]

\medskip

\[
\begin{bmatrix}
1 & 2 & -1\\
0 & -3 & 2\\
0 & 0 & 3
\end{bmatrix}
\begin{bmatrix}
x_1 \\ x_2 \\ x_3
\end{bmatrix}
=
\begin{bmatrix}
0 \\ 0 \\ 1
\end{bmatrix}
\]
\[
\implies
x_3 = \frac13, \; x_2 = \frac{0-2(\frac13)}{-3} = \frac29, \; x_1 = \frac{0+1(\frac13)-2(\frac29)}{1}=-\frac19.
\]

\medskip
\(
 \therefore A^{-1}=\begin{bmatrix}
 -\dfrac29&\phantom{-}\dfrac59&-\dfrac19\\[3mm]
 \phantom{-}\dfrac49&-\dfrac19&\phantom{-}\dfrac29\\[3mm]
 -\dfrac13&\phantom{-}\dfrac13&\phantom{-}\dfrac13
 \end{bmatrix}.
\)
\end{solution}

\subsection{Matrix Factorisation}
Gaussian elimination is the principal tool in the direct solution of linear system of equations. Previously, we have learnt that the steps used for elimination can also be used to factor a matrix.
Suppose the matrix \(A\) can be factorized as $A=LU$,
where $L$ is lower triangular and $U$ is upper triangular.  Then
\[
 A\vect{x}=\vect b
 \quad\Longleftrightarrow\quad
 (LU)\vect{x}=\vect b
 \quad\Longleftrightarrow\quad
 L(U\vect{x})=\vect b.
\]
Let $U\vect{x}=\vect y$.  Now the problem $A\vect x = \vect b$ has been separated into two equations $L\vect y = \vect b \text{ and } U\vect x = \vect y$. Since both $L$ and $U$ are triangular, both equations could be solved by substitution operations (forward and backward substitution respectively).

Since substitution has the operation count $\sim O(n^2)$, it is expected to speed up the process of solving $A\vect x = \vect b$. However, for only one value for $\vect b$, matrix factorization does not have much advantage over naive Gaussian Elimination. But, for different values of $\vect{b}_i$, matrix factorization would be preferred as one only needs to factorize $A$ once and the same $L$ and $U$ could be applied to the different values of $\vect{b}_i$ repeatedly.
\begin{example}
Factorise
\[
 A=\begin{bmatrix}
 1&1&0&3\\
 2&1&-1&1\\
 3&-1&-1&2\\
 -1&2&3&-1
 \end{bmatrix}
\]
into $LU$ form and solve $A\vect{x}=\vect b$ for
$\vect b=[1,1,-3,4]^T$.
\end{example}

\begin{solution}[8.0cm]
\[
A=\begin{bmatrix}
1 & 1 & 0 & 3\\
2 & 1 & -1 & 1\\
3 & -1 & -1 & 2\\
-1 & 2 & 3 & -1
\end{bmatrix}
\xrightarrow[
{\displaystyle
\begin{array}{l}
R_4+R_1\to R_4
\end{array}}]
{{\displaystyle
\begin{array}{l}
R_2-2R_1\to R_2\\
R_3-3R_1\to R_3
\end{array}}
}
\begin{bmatrix}
1 & 1 & 0 & 3\\
0 & -1 & -1 & -5\\
0 & -4 & -1 & -7\\
0 & 3 & 3 & 2
\end{bmatrix}
\xrightarrow[
{\displaystyle
\begin{array}{l}
R_4+3R_2\to R_4
\end{array}}]
{\displaystyle
\begin{array}{l}
R_3-4R_2\to R_3
\end{array}}
\]
\[
\begin{bmatrix}
1 & 1 & 0 & 3\\
0 & -1 & -1 & -5\\
0 & 0 & 3 & 13\\
0 & 0 & 0 & -13
\end{bmatrix}
=U,
\qquad
L=\begin{bmatrix}
1 & 0 & 0 & 0\\
2 & 1 & 0 & 0\\
3 & 4 & 1 & 0\\
-1 & -3 & 0 &1
\end{bmatrix}
\]
\(\therefore A\vect x = \vect b \implies LU\vect x = \vect b.\)

\noindent
Let $\vect y = U\vect x$. Then $LU\vect x = \vect b \implies L\vect y = \vect b$.
\[
\implies
\begin{bmatrix}
1 & 0 & 0 & 0\\
2 & 1 & 0 & 0\\
3 & 4 & 1 & 0\\
-1 & -3 & 0 & 1\\
\end{bmatrix}
\begin{bmatrix}
y_1 \\ y_2 \\ y_3 \\ y_4
\end{bmatrix}
=
\begin{bmatrix}
1 \\ 1 \\ -3 \\ 4
\end{bmatrix}
\implies
 y_1=1,
 \quad y_2=1-2(1)=-1,
\]
\[
 \quad y_3=-3-3(1)-4(-1)=-2,
 \quad y_4=4+1+3(-1)+0=2.
\]
\(
\implies U\vect x = \vect y \implies
\begin{bmatrix}
1 & 1 & 0 & 3\\
0 & -1 & -1 & -5\\
0 & 0 & 3 & 13\\
0 & 0 & 0 & -13\\
\end{bmatrix}
\begin{bmatrix}
x_1 \\ x_2 \\ x_3 \\ x_4
\end{bmatrix}
=
\begin{bmatrix}
1 \\ -1 \\ -2 \\ 2
\end{bmatrix}
\)

\smallskip
\noindent
\(\displaystyle
\implies
 x_4=-\frac2{13},
 \quad x_3=-2-13\left(-\frac{2}{13}\right)=0,
 \quad x_2=\frac1{-1}\left(-1+5\left(-\frac2{13}\right)+0\right)=\frac{23}{13},
\)

\smallskip
\noindent
\(\displaystyle
\; \; \qquad x_1=1-3\left(-\frac2{13}\right)-0-\frac{23}{13}=-\frac4{13}.
\)

\[
 \therefore \vect x=\left[-\frac4{13},\frac{23}{13},0,-\frac2{13}\right]^T.
\]
\end{solution}

\section{Iterative Methods for Solving Linear Systems}

In this section, iterative methods start with an initial guess or approximation to a solution, $\vect{x}^{(0)}$ and then generate a succession of better and better approximations
\[
 \vect{x}^{(1)},\vect{x}^{(2)},\vect{x}^{(3)},\ldots
\]
that (may) tend toward the exact solution.
We will discuss three iterative methods:
\begin{enumerate}
\item Jacobi method
\item Gauss-Seidel method
\item Successive Over-relaxation (SOR) method
\end{enumerate}

\noindent
To get the idea, we separate our coefficient matrix A from our linear system as the following:
\[
 A=D+L+U,
\]
where $D$ contains the diagonal entries of $A$, $L$ the entries below the diagonal, and $U$ the entries above the diagonal.

\noindent
Don't worry if the formulas found in the explanations below feel a bit overwhelming. For first time learners, referring to the example to get a feel how the methods work will clear up a lot of the confusion. The same goes for the majority of sections that follow in future chapters. The formulas are just there to properly define how the methods work in mathematical terms (that may be useful when trying to implement the methods in programming).

\subsection{Jacobi Method}
Consider a linear system with $n$ equations and $n$ unknowns. Then we first solve for $x_i$ in the $i$-th equation (where $i=1,2,\dots,n$) to derive our iteration equation.
Componentwise,
\[
 x_i^{(k+1)}
 =\frac{1}{a_{ii}}
 \left(b_i-\sum_{j\ne i}a_{ij}x_j^{(k)}\right).
\]
In matrix notation, the Jacobi iteration is
\[
 \vect{x}^{(k+1)}
 =D^{-1}\bigl(\vect b-(L+U)\vect{x}^{(k)}\bigr).
\]
The order in which the equations are examined is irrelevant, since the Jacobi method treats them independently. That is, we prepare the vector $\vect x^{(k+1)}$ only using $\vect x^{(k)}$, or iterating in batches, only with the information from the previous iteration. For this reason, the Jacobi method is also known as the \textbf{method of simultaneous corrections}, since the updates could in principle be done simultaneously.

\subsection{Gauss--Seidel Method}
Similar to the Jacobi method, we prepare the iteration equations by solving for $x_i$ in the $i$-th equation of our linear system. Componentwise,
\[
 x_i^{(k+1)}=\frac{1}{a_{ii}}
 \left(
 b_i-\sum_{j<i}a_{ij}x_j^{(k+1)}
     -\sum_{j>i}a_{ij}x_j^{(k)}
 \right).
\]
In matrix notation, the Gauss--Seidel iteration is
\[
 (D+L)\vect{x}^{(k+1)}
 =\vect b-U\vect{x}^{(k)}.
\]
Unlike the Jacobi method where $\vect x^{(k)}$ is prepared in batches, each element of $\vect x^{(k)}$, $x^{(k)}_i$ depends on $x^{(k)}_j \; \forall j<i$; that is we iterate with the \textbf{latest information} instead of information from the previous iteration. 
This is a \textbf{method of successive corrections}; it \textbf{replaces approximations by corresponding new ones} as soon as the latter are available.
Note that since newly computed values are used immediately, so Gauss--Seidel often converges faster than Jacobi.

\subsection{Successive Over-Relaxation}
Modification of the Gauss-Seidel method with the relaxation parameter $\omega \; (0<\omega<2)$ is done to improve the rate of convergence. Componentwise, we only add weights as follows:
\[
 x_i^{(k+1)}
 =
 \frac{\omega}{a_{ii}}
 \left(
 b_i-\sum_{j<i}a_{ij}x_j^{(k+1)}
     -\sum_{j>i}a_{ij}x_j^{(k)}
 \right)+(1-\omega)x_i^{(k)},
\qquad 0<\omega<2.
\]
\begin{itemize}
\item If $\omega=1$, SOR reduces to Gauss--Seidel.
\item Values $0<\omega<1$ under-relax the iteration, which might help to dampen the iteration to achieve convergence if the Gauss-Seidel method diverges.
\item Meanwhile, if the Gauss-Seidel already converges, values $1<\omega<2$ will help to speed up the convergence by over-relaxation.
\end{itemize}

\begin{definition}[Strict diagonal dominance]
A square matrix $A=[a_{ij}]$ of order $n$ is strictly diagonally dominant if
\[
 |a_{kk}| >\sum_{\substack{i=1\\ i\ne k}}^{n}|a_{ki}|,
 \qquad k=1,2,\ldots,n.
\]
That is, A is strictly diagonally dominant if the absolute value of each diagonal entry is greater than the sum of the absolute values of the remaining entries in the same row.
\end{definition}

\begin{example}
Show that
\(
 A=\begin{bmatrix}
 2&-7&4\\
 8&1&6\\
 -3&5&9
 \end{bmatrix}
\)
is not strictly diagonally dominant.
\end{example}

\begin{solution}[2.5cm]
We see from the first and second row that
\(
 |2|<|-7|+|4|=11
\)
and $|1|<|8|+|6|=14$. 
Hence, $A$ is not strictly diagonally dominant.
\end{solution}

\begin{theorem}[Convergence of the Iterative Methods]
If a square matrix $A$ is strictly diagonally dominant, then the Jacobi and Gauss--Seidel approximations to the solution of the linear system $A\vect x = \vect b$ both converge to the exact solution for all choices of the initial approximation or guess, $\vect{x}^{(0)}$.
\end{theorem}

\noindent
Before we finally demonstrate how to execute these iterative methods in practice, we need to define a termination criterion. For these methods in this course, we commonly stop the computation when
\[
    \max_i|x_i^{(k+1)}-x_i^{(k)}|<\epsilon, 
    \quad \text{where } \epsilon \text{ is the pre-specified stopping criterion}.
\]
In practice however, a residual test such as $\norm{\vect b-A\vect{x}^{(k)}}<\epsilon$ may also be considered because it measures how well the computed vector satisfies the original equations.

\begin{example}
Use the \textbf{Jacobi method}, to solve
\[
\begin{aligned}
20x_1+x_2-x_3&=17\\
x_1-10x_2+x_3&=13\\
-x_1+x_2+10x_3&=18
\end{aligned}
\]
Iterate until $|x_i^{[p+1]}-x_i^{[p]}| < 0.0002$, for all $i$. Prepare all the computations in 5 decimal places.
\end{example}

\begin{solution}[8.0cm]
\[
 x_1^{[k+1]}=\frac{17-x_2^{[k]}+x_3^{[k]}}{20},
 \quad
 x_2^{[k+1]}=\frac{13-x_1^{[k]}-x_3^{[k]}}{-10},
 \quad
 x_3^{[k+1]}=\frac{18+x_1^{[k]}-x_2^{[k]}}{10}.
\]
Let our initial guess be: $x_1^{[0]}=x_2^{[0]}=x_3^{[0]}=0$. Then

\smallskip
\(\displaystyle
x_1^{[1]}=\frac{17}{20}=0.85000, \quad x_2^{[1]}=-\frac{13}{10}=-1.30000, \quad x_3^{[1]}=\frac{18}{10}=1.80000.
\)

\smallskip
\(
|x_1^{[1]} - x_1^{[0]}| = |0.85000-0| = 0.85 \geq 0.0002.
\)

\bigskip
\(\displaystyle
x_1^{[2]}=\frac{17-(-1.3)+1.8}{20}=1.00500, \quad x_2^{[2]}=-\frac{13-0.85-1.8}{10}=-1.03500,\)

\smallskip
\(\displaystyle x_3^{[2]}=\frac{18+0.85-(-1.3)}{10}=2.01500.
\)

\smallskip\(
|x_1^{[2]} - x_1^{[1]}| = |1.00500-0.85000| = 0.15500 \geq 0.0002.
\)

\bigskip
\(\displaystyle
x_1^{[3]}=\frac{17-(-1.035)+2.015}{20}=1.00250, \quad x_2^{[3]}=-\frac{13-1.005-2.015}{10}=-0.99800,\)

\smallskip
\(\displaystyle x_3^{[3]}=\frac{18+1.005-(-1.035)}{10}=2.00400.
\)

\smallskip\(
|x_1^{[3]} - x_1^{[2]}| = |1.00250-1.00500| = 0.0025 \geq 0.0002.
\)

\bigskip
\(\displaystyle
x_1^{[4]}=\frac{17-(-0.998)+2.004}{20}=1.00010, \quad x_2^{[4]}=-\frac{13-1.0025-2.004}{10}=-0.99935,\)

\smallskip
\(\displaystyle x_3^{[4]}=\frac{18+1.0025-(-0.998)}{10}=2.00005.
\)

\smallskip\(
|x_1^{[4]} - x_1^{[3]}| = |1.00010-1.00250| = 0.00240 \geq 0.0002.
\)

\bigskip
\(\displaystyle
x_1^{[5]}=\frac{17-(-0.99935)+2.00005}{20}=0.99997, \)

\smallskip
\(\displaystyle x_2^{[5]}=-\frac{13-1.0001-2.00005}{10}=-0.99999,\)

\smallskip
\(\displaystyle x_3^{[5]}=\frac{18+1.00010-(-0.99935)}{10}=1.99994.
\)

\smallskip\(
|x_1^{[5]} - x_1^{[4]}| = |0.99997-1.00010| = 0.00013 < 0.0002.
\)

\smallskip\(
|x_2^{[5]} - x_2^{[4]}| = |-0.99998-(-0.99935)| = 0.00063 \geq 0.0002.
\)

\bigskip
\(\displaystyle
x_1^{[6]}=\frac{17-(-0.99999)+1.99994}{20}=1.00000, \)

\smallskip
\(\displaystyle x_2^{[6]}=-\frac{13-0.99997-1.99994}{10}=-1.00000,\)

\smallskip
\(\displaystyle x_3^{[6]}=\frac{18+0.99997-(-0.99999)}{10}=2.00000.
\)

\smallskip\(
|x_1^{[6]} - x_1^{[5]}| = |1.00000-0.99997| = 0.00003 < 0.0002.
\)

\smallskip\(
|x_2^{[6]} - x_2^{[5]}| = |-1-(-0.99998)| = 0.00002 < 0.0002.
\)

\smallskip\(
|x_3^{[6]} - x_3^{[5]}| = |2.00000-1.99994| = 0.00006 < 0.0002.
\)

\medskip
\noindent
Hence, $x_1 = 1, x_2 = -1, x_3 =2$.
\end{solution}

\begin{example}
Redo the previous example but by using the Gauss--Seidel method.
\end{example}

\begin{solution}[6.5cm]
\[
 x_1^{[k+1]}=\frac{17-x_2^{[k]}+x_3^{[k]}}{20},
 \quad
 x_2^{[k+1]}=\frac{13-x_1^{[k+1]}-x_3^{[k]}}{-10},
 \quad
 x_3^{[k+1]}=\frac{18+x_1^{[k+1]}-x_2^{[k+1]}}{10}.
\]
Let our initial guess be: $x_1^{[0]}=x_2^{[0]}=x_3^{[0]}=0$. Then

\smallskip
\(\displaystyle
x_1^{[1]}=\frac{17}{20}=0.85000, \quad x_2^{[1]}=-\frac{13-0.85}{10}=-1.21500, \)

\smallskip
\(\displaystyle x_3^{[1]}=\frac{18+0.85-(-1.215)}{10}=2.00650.\)

\smallskip
\(
|x_1^{[1]} - x_1^{[0]}| = |0.85000-0| = 0.85 \geq 0.0002.
\)

\bigskip
\(\displaystyle
x_1^{[2]}=\frac{17-(-1.215)+2.0065}{20}=1.01108, \quad x_2^{[2]}=-\frac{13-1.01108-2.0065}{10}=-0.99824,\)

\smallskip
\(\displaystyle x_3^{[2]}=\frac{18+1.01108-(-0.99824)}{10}=2.00093.
\)

\smallskip\(
|x_1^{[2]} - x_1^{[1]}| = |1.01108-0.85000| = 0.16108 \geq 0.0002.
\)

\bigskip
\(\displaystyle
x_1^{[3]}=\frac{17-(-0.99824)+2.00093}{20}=0.99996, \)

\smallskip
\(\displaystyle x_2^{[3]}=-\frac{13-0.99996-2.00093}{10}=-0.99991,\)

\smallskip
\(\displaystyle x_3^{[3]}=\frac{18+0.99996-(-0.99991)}{10}=1.99999
\)

\smallskip\(
|x_1^{[3]} - x_1^{[2]}| = |1.01108-0.99996| = 0.01112 \geq 0.0002.
\)

\bigskip
\(\displaystyle
x_1^{[4]}=\frac{17-(-0.99991)+1.99999}{20}=1.00000, \quad x_2^{[4]}=-\frac{13-1-1.99999}{10}=-1.00000,\)

\smallskip
\(\displaystyle x_3^{[4]}=\frac{18+1-(-1)}{10}=2.00000.
\)

\smallskip\(
|x_1^{[4]} - x_1^{[3]}| = |1.00000-0.99996| = 0.00004 < 0.0002.
\)

\smallskip\(
|x_2^{[4]} - x_2^{[3]}| = |-1.00000-(-0.99991)| = 0.00009 < 0.0002.
\)

\smallskip\(
|x_3^{[4]} - x_3^{[3]}| = |2.00000-1.99999| = 0.00001 < 0.0002.
\)

%\[
%\begin{array}{c|rrr}
%\toprule
%k&x_1^{(k)}&x_2^{(k)}&x_3^{(k)}\\
%\midrule
%0&0&0&0\\
%1&0.85000&-1.21500&2.00650\\
%\midrule
%2&1.01108&-0.99824&2.00093\\
%3&0.99996&-0.99991&1.99999\\
%4&0.99999&-1.00000&2.00000\\
%\bottomrule
%\end{array}
%\]
\medskip \noindent
Hence, $x_1 = 1, x_2 = -1, x_3 = 2$.
\end{solution}

\begin{example}
Use two iterations of \textbf{Successive Over-relaxation (SOR) method} to solve
\[
\begin{aligned}
3x_1-x_2+x_3&=-1,\\
-x_1+3x_2-x_3&=7,\\
x_1-x_2+3x_3&=-7.
\end{aligned}
\]
with $\omega = 1.25$ and initial approximation $x_1^{[0]}=0,x_2^{[0]}=0,x_3^{[0]}=0$.
\end{example}

\begin{solution}[5.5cm]
With $\omega=1.25$, we have:

\smallskip
\noindent
\(\displaystyle
x_1^{[k+1]}=\frac{1.25}{3}(-1+x_2^{[k]}-x_3^{[k]})-0.25x_1^{[k]}, 
\quad x_2^{[k+1]}=\frac{1.25}{3}(7+x_1^{[k+1]}+x_3^{[k]})-0.25x_2^{[k]},
\)

\smallskip \noindent
\(\displaystyle x_3^{[k+1]}=\frac{1.25}{3}(-7-x_1^{[k+1]}+x_2^{[k+1]})-0.25x_3^{[k]}.\)

\noindent Then

\smallskip
\(\displaystyle
x_1^{[1]}=\frac{1.25}{3}(-1)-0=-0.4167, \quad x_2^{[1]}=\frac{1.25}{3}(7-0.4167)-0=2.7431,
\)

\smallskip
\(\displaystyle
x_3^{[1]}=\frac{1.25}{3}(-7-(-0.4167)+2.7431)-0=-1.6001.
\)

\bigskip
\(\displaystyle
x_1^{[2]}=\frac{1.25}{3}(-1+2.7431-(-1.6001))-0.25(-0.4167)=1.4972, \)

\smallskip
\(\displaystyle x_2^{[2]}=\frac{1.25}{3}(7+1.4972-1.6001)-0.25(2.7431)=2.1880,
\)

\smallskip
\(\displaystyle
x_1^{[3]}=\frac{1.25}{3}(-7-1.4972+2.1880)-0.25(-1.6001)=-2.2288.
\)

\medskip \noindent
Hence, $x_1 = 1.4972, x_2 = 2.188, x_3 = -2.2288$.
\end{solution}

\subsection{Conjugate Gradient Method}
The Conjugate Gradient Method is a popular iterative method for solving a sparse system. Consider a quadratic form and its gradient:
\[
f(\vect x) = \frac12 \vect x^TA \vect x - \vect x^T \vect b,
\qquad \nabla f(\vect x) = \frac12 (A^T\vect x + A\vect x) - \vect b.
\]
If $A$ is symmetric, then $\nabla f(\vect x) = A\vect x - \vect b$.
Therefore, $f(\vect x)$ has a critical point at $\vect x$ if $A\vect x = \vect b$. Thus, if $A$ is a positive, definite, and symmetric, then we can solve $A\vect x = \vect b$ by minimising $f(x)$.

In Machine Learning, the method of steepest descent (or more widely known as gradient descent) is an another algorithm used to minimise a loss function $f(\vect x)$. 
But it converges much more slowly towards to exact solution.
On the other hand, the conjugate gradient method is able to reach the \textbf{exact solution} in at most $\mathbf{n}$ steps (given $A$ is a square matrix of order $n$).
\textbf{This section will be rewritten in the future due to having too many gaps in understanding.}
%On the other hand, the conjugate gradient method constructs mutually $A$-conjugate search directions and, in exact arithmetic, reaches the exact solution in at most $n$ steps.

\begin{algorithm}[Conjugate gradient]
Start with an initial approximation $\vect{x}_0$. Then
\[
 1.\;\vect r_0=\vect b-A\vect{x}_0,
 \qquad
 2.\;\mathbf{P}_0=\vect r_0.
\]
For $k=0,1,2,\ldots,n$, compute
\[
 \alpha_k=\frac{\vect r_k\T\vect r_k}{\vect p_k\T A\vect p_k},
 \qquad
 \vect{x}_{k+1}=\vect{x}_k+\alpha_k\vect p_k,
\]
\[
 \vect r_{k+1}=\vect r_k-\alpha_kA\vect p_k,
 \qquad
 \beta_k=\frac{\vect r_{k+1}\T\vect r_{k+1}}{\vect r_k\T\vect r_k},
\]
\[
 \vect p_{k+1}=\vect r_{k+1}+\beta_k\vect p_k.
\]
\end{algorithm}

\begin{example}
Apply conjugate gradient with $\vect{x}_0=\vect0$ to
\[
 A=\begin{pmatrix}
4&3&0\\3&4&-1\\0&-1&4
\end{pmatrix},
\qquad
\vect b=\begin{pmatrix}24\\30\\-24\end{pmatrix}.
\]
\end{example}

\begin{solution}[7.0cm]
The successive approximations are
\[
\begin{array}{c|rrr}
\toprule
k&x_1&x_2&x_3\\
\midrule
0&0&0&0\\
1&3.52577&4.40722&-3.52577\\
2&2.85801&4.14897&-4.95422\\
3&3.00000&4.00000&-5.00000\\
\bottomrule
\end{array}
\]
The residual after the third step is zero to round-off accuracy, giving
\[
 \vect{x}=(3,4,-5)\T.
\]
\end{solution}

\section{Eigenvalues and Eigenvectors}

\begin{definition}[Eigenvalues and Eigenvectors]
A scalar $\lambda$ is an eigenvalue of an $n\times n$ matrix $A$ if there exists a nonzero vector $\vect v \in \mathbb{R}^n$ such that
\[
 A\vect v=\lambda\vect v.
\]
The nonzero vector $\vect v$ is an eigenvector corresponding to $\lambda$.
\end{definition}

\noindent
Note that eigenvalues satisfy the characteristic equation
\(
 \det(A-\lambda I)=0.
\)
Once $\lambda$ is known, its corresponding eigenvectors are the nonzero solutions of
$(A-\lambda I)\vect v=\vect0$.

\begin{example}
For
\[
 A=\begin{bmatrix}4&1\\3&2\end{bmatrix},
 \qquad
 \vect v=\begin{bmatrix}1\\1\end{bmatrix},
\]
verify that $\vect v$ is an eigenvector of $A$ and find its corresponding eigenvalue.
\end{example}

\begin{solution}[2.2cm]
\[
 A\vect v
 =\begin{bmatrix}4&1\\3&2\end{bmatrix}\begin{bmatrix}1\\1\end{bmatrix}
 =\begin{bmatrix}5\\5\end{bmatrix}
 =5\begin{bmatrix}1\\1\end{bmatrix}.
\]
Hence $\vect v$ is an eigenvector corresponding to $\lambda=5$.
\end{solution}

\subsection{Gerschgorin Discs}

\begin{theorem}[Gercshgorin's theorem]
The eigenvalues of an $n\times n$ matrix $A$ lie in the union of $n$ discs
\[
 \abs{\lambda-a_{ii}}\le r_i,
 \quad
 r_i=\sum_{j\ne i}|a_{ij}| = |a_{i1}| + \dots + |a_{i,i-1}| + |a_{i,i+1}| + \dots + |a_{in}|.
\]
\end{theorem}

Gerschgorin's Theorem (or Gershgorin's Theorem) is a finding that is common taught in Complex Analysis that is used to estimate the location of the eigenvalues of a matrix $A$. For a general square matrix, the union of the $n$ disks would be drawn as circles on the complex plane. However, for a real symmetric matrix, all eigenvalues are real, so the discs may be viewed as intervals on the real line.

\begin{example}
Find Gerschgorin intervals corresponding to
\[
 A=\begin{bmatrix}
7&4&-4\\
4&-8&-1\\
-4&-1&-8
\end{bmatrix}.
\]
\end{example}

\begin{solution}[3.8cm]
The row radii are
\[
 r_1=|4|+|-4|=8,
 \qquad r_2=|4|+|-1|=5,
 \qquad r_3=|-4|+|-1|=5.
\]
Thus we have
\[
|\lambda_1 - 7| \leq 8 \implies -8+7 \leq \lambda_1 \leq 8+1 \implies -1 \leq \lambda_1 \leq 15.
\]
\[
|\lambda_2 - (-8)| \leq 5 \implies -5-8 \leq \lambda_1 \leq -5+8 \implies -13 \leq \lambda_2 \leq -3.
\]
\[
|\lambda_3 - (-8)| \leq 5 \implies -5-8 \leq \lambda_3 \leq -5+8 \implies -13 \leq \lambda_3 \leq -3.
\]

\(\displaystyle \therefore \text{The eigenvalues of A, } \lambda \in [-13, -3] \cup [-1,15]\).
\end{solution}

\subsection{Power Method}

\begin{definition}[Dominant eigenvalue]
An eigenvalue $\lambda_1$ is dominant if
\[
 |\lambda_1|>|\lambda_j|
 \qquad\text{for every }j\ne1.
\]
That is if the absolute value of $\lambda$ is larger than the absolute values of all the other remaining eigenvalues. An eigenvector corresponding to the dominant eigenvalue is thus also called a dominant eigenvector.
\end{definition}

\begin{example}
We explore the concept of dominant eigenvalues in the following cases:
\begin{enumerate}[label=(\alph*),leftmargin=2.7em]
\item A $4\times 4$ matrix A with eigenvalues $\lambda=-5,4,2,-3$ has dominant eigenvalue $\lambda = -5$ since 
\[\lvert -5\rvert > \lvert 4\rvert,\quad
\lvert -5\rvert > \lvert 2\rvert,\quad
\lvert -5\rvert > \lvert -3\rvert.\]
\item A $3\times 3$ matrix $A$ with eigenvalues $\lambda=2,-5,5$ has no dominant eigenvalue since \[\lvert -5 \rvert \ngtr \lvert 5 \rvert.\]
\end{enumerate}
\end{example}

\begin{theorem}[The Power Method or Direct Iteration Method]
If an $n\times n$ matrix $A$ is diagonalizable (i.e. having $n$ distinct eigenvalues) and the initial guess for the dominant eigenvector $\vect v$, $\vect x_0$ has a nonzero component in the dominant eigendirection, then
\[
\vect x_k = A^k\vect{x}_0 = A \vect x_{k-1}
\]
is a good approximation to a dominant eigenvector when $k$ is sufficiently large.
\end{theorem}

\begin{enumerate}[label=(\alph*)]
\item The dominant eigenvalue may be approximated using the current approximation to the dominant eigenvector $\vect x_k = A^k\vect x_0$ by the \textbf{Rayleigh quotient}
\[
\lambda \approx \frac{\vect{x}_k^T A\vect{x}_k}{\vect{x}_k^T\vect{x}_k} = \frac{A\vect x_k \cdot \vect x_k}{\vect x_k \cdot \vect x_k}.
\]
\item The \textbf{convergence rate} depends on the ratios of the other remaining eigenvectors to the dominant eigenvector
\[
\frac{\lambda_2}{\lambda}, \dots , \frac{\lambda_n}{\lambda}.
\]
The smaller the ratios, the faster the rate of convergence.
\end{enumerate}


\begin{example}
Let
\[
 A=\begin{bmatrix}4&-2\\3&-1\end{bmatrix},
 \qquad
 \vect{x}_0=\begin{bmatrix}1\\0\end{bmatrix}.
\]
Use the power method to approximate the dominant eigenvalue and a corresponding eigenvector with the initial guess given above. Prepare 6 iterations.
\end{example}

\begin{solution}[6.8cm]
\(\vect{x}_1=A\vect x_0 = \begin{bmatrix}4&-2\\3&-1\end{bmatrix}
\begin{bmatrix}1\\0\end{bmatrix}=\begin{bmatrix}4\\3\end{bmatrix}\)

\smallskip \noindent
\(\vect{x}_2=\begin{bmatrix}4&-2\\3&-1\end{bmatrix}\begin{bmatrix}4\\3\end{bmatrix}
=\begin{bmatrix}10\\9\end{bmatrix}\)

\smallskip \noindent
\(\vect{x}_3=\begin{bmatrix}4&-2\\3&-1\end{bmatrix}\begin{bmatrix}10\\9\end{bmatrix}
=\begin{bmatrix}22\\21\end{bmatrix}\)

\smallskip \noindent
\(\vect{x}_4=\begin{bmatrix}4&-2\\3&-1\end{bmatrix}\begin{bmatrix}22\\21\end{bmatrix}
=\begin{bmatrix}46\\45\end{bmatrix}\)

\smallskip \noindent
\(\vect{x}_5=\begin{bmatrix}4&-2\\3&-1\end{bmatrix}\begin{bmatrix}46\\45\end{bmatrix}
=\begin{bmatrix}94\\93\end{bmatrix}\)

\smallskip \noindent
\(\vect{x}_6=\begin{bmatrix}4&-2\\3&-1\end{bmatrix}\begin{bmatrix}94\\93\end{bmatrix}
=\begin{bmatrix}190\\189\end{bmatrix}\)

\begin{flalign*}
\therefore \text{The dominant eigenvalue, }\lambda
&\approx
\frac{\mathbf{x}_6^T A\mathbf{x}_6}
     {\mathbf{x}_6^T\mathbf{x}_6}=
\frac{
\begin{bmatrix}190\\189\end{bmatrix}^{T}
\begin{bmatrix}4&-2\\3&-1\end{bmatrix}
\begin{bmatrix}190\\189\end{bmatrix}
}{
\begin{bmatrix}190\\189\end{bmatrix}^{T}
\begin{bmatrix}190\\189\end{bmatrix}
}
=
\frac{
\begin{bmatrix}190&189\end{bmatrix}
\begin{bmatrix}382\\381\end{bmatrix}
}{
190^2+189^2
} &&\\[4pt]
&=
\frac{190(382)+189(381)}
     {190^2+189^2}
=2.0132. &&
\end{flalign*}
\end{solution}

\noindent
We can observe some key facts from the previous example:
\begin{enumerate}
\item It is clear that the vectors $x_m$ are getting closer and closer to scalar multiples of
    $\begin{bmatrix}1\\1\end{bmatrix}$ which is a dominant eigenvector of $A$. By checking, we also see that the corresponding dominant eigenvalue, $\lambda = 2$.
\item The convergence is rather slow.
\item The power method often generates a sequence of vectors $\{A^k\vect x_0\}$ that have inconveniently large entries.
\end{enumerate}

\begin{theorem}[Power Method with Scaling]
The large entries discussed earlier can be avoided by scaling the iterative vector at each step. 

\smallskip
\noindent
That is, we first compute
\(\displaystyle
\vect y_{k} = A\vect{x}_{k-1} = [y_1, \dots, y_n]^T \text{ and multiply it by }
\frac1{ \max \{|y_{1}|,\dots,|y_{n}|\}}\) to get
\[
    \vect x_{k+1} = \frac1{\max\{|A\vect x_k|\}}A\vect x_{k}.
\]
\end{theorem}
We repeat this process to obtain scaled-down vectors, and the approximate dominant eigenvalue (at step $k$) is computed using the Rayleigh quotient:
\[
\lambda (k) = \frac{\vect x_k^T A \vect x_k}{\vect x_k^T \vect x_k} = \frac{A\vect x_k \cdot \vect x_k}{\vect x_k \cdot \vect x_k}.
\]
We will \textbf{STOP} the computation at the $k$-th step if the \textbf{estimated relative error} or \textbf{estimated percentage error} is low enough:
\[
\left| \frac{\lambda(k)-\lambda(k-1)}{\lambda(k)} \right| < \epsilon, \quad
\text{ or } \quad \left| \frac{\lambda(k)-\lambda(k-1)}{\lambda(k)} \right| < \epsilon, \quad
\text{ for a small } \epsilon.
\]

\begin{example}
Repeat the previous example using power iteration with scaling.
\end{example}

\begin{solution}[3.6cm]
\(\displaystyle
A\vect x_0 = \begin{bmatrix}4&-2\\3&-1\end{bmatrix}
\begin{bmatrix}1\\0\end{bmatrix}=\begin{bmatrix}4\\3\end{bmatrix},
\quad \vect x_1 = \frac14 \begin{bmatrix}4\\3\end{bmatrix}=\begin{bmatrix}1\\0.75\end{bmatrix}\)

\smallskip \noindent
\(\displaystyle
A\vect x_1=\begin{bmatrix}4&-2\\3&-1\end{bmatrix}\begin{bmatrix}1\\0.75\end{bmatrix}
=\begin{bmatrix}2.5\\2.25\end{bmatrix},
\quad \vect x_2 = \frac1{2.5}\begin{bmatrix}2.5\\2.25\end{bmatrix} = \begin{bmatrix}1\\0.9\end{bmatrix}\)

\smallskip \noindent
\(\displaystyle
A\vect x_2=\begin{bmatrix}4&-2\\3&-1\end{bmatrix}\begin{bmatrix}1\\0.9\end{bmatrix}
=\begin{bmatrix}2.2\\2.1\end{bmatrix},
\quad \vect x_3 = \frac1{2.2}\begin{bmatrix}2.2\\2.1\end{bmatrix} = \begin{bmatrix}1\\0.9545\end{bmatrix}\)

\smallskip \noindent
\(\displaystyle
A\vect x_3=\begin{bmatrix}4&-2\\3&-1\end{bmatrix}\begin{bmatrix}1\\0.9545\end{bmatrix}
=\begin{bmatrix}2.0910\\2.0455\end{bmatrix},
\quad \vect x_4 = \frac1{2.0910}\begin{bmatrix}2.0910\\2.0455\end{bmatrix} = \begin{bmatrix}1\\0.9782\end{bmatrix}\)

\smallskip \noindent
\(\displaystyle
A\vect x_4=\begin{bmatrix}4&-2\\3&-1\end{bmatrix}\begin{bmatrix}1\\0.9782\end{bmatrix}
=\begin{bmatrix}2.0436\\2.0218\end{bmatrix},
\quad \vect x_5 = \frac1{2.0436}\begin{bmatrix}2.0436\\2.0218\end{bmatrix} = \begin{bmatrix}1\\0.9893\end{bmatrix}\)

\smallskip \noindent
\(\displaystyle
A\vect x_5=\begin{bmatrix}4&-2\\3&-1\end{bmatrix}\begin{bmatrix}1\\0.9893\end{bmatrix}
=\begin{bmatrix}2.0214\\2.0107\end{bmatrix},
\quad \vect x_6 = \frac1{2.0214}\begin{bmatrix}2.0214\\2.0107\end{bmatrix} = \begin{bmatrix}1\\0.9947\end{bmatrix}\)

\begin{flalign*}
\lambda(6)&=\frac{A\vect x_6 \cdot \vect x_6}{\vect x_6 \cdot \vect x_6}
=\frac{\begin{bmatrix}4 & -2\\3 & -1\end{bmatrix}\begin{bmatrix}1\\0.9947\end{bmatrix}\cdot\begin{bmatrix}1\\0.9947\end{bmatrix}}{\begin{bmatrix}1\\0.9947\end{bmatrix}\cdot\begin{bmatrix}1\\0.9947\end{bmatrix}}
=\frac{\begin{bmatrix}2.0106\\2.0053\end{bmatrix}\cdot\begin{bmatrix}1\\0.9947\end{bmatrix}}{1^2+0.9947^2}\\
&=\frac{2.0106 \cdot 1 + 2.0053 \cdot 0.9947}{1^2 + 0.9947^2}=2.0133.
\end{flalign*}

Hence, using the Power Method with Scaling, we obtain the dominant eigenvalue and an eigenvector of $A$ corresponding to the dominant eigenvalue:
\[
\lambda \approx 2.0133
\qquad \vect x \approx \begin{bmatrix}1\\0.9947\end{bmatrix}.
\]
\end{solution}

\subsection{Inverse and Shifted Inverse Power Methods}

If $A$ is nonsingular, the eigenvalues of $A^{-1}$ are $1/\lambda_i$.  Hence power iteration applied to $A^{-1}$ finds the eigenvalue of $A$ having smallest modulus. In practice one solves
\[
 A\vect y_k=\vect{x}_{k-1}
\]
rather than explicitly forming $A^{-1}$.

\begin{example}
Let
\[
 A=\begin{pmatrix}2&3\\4&1\end{pmatrix},
 \qquad
 \vect{x}_0=(-0.7,1)\T.
\]
Use two inverse power iterations with scaling to approximate the smaller eigenvalue.
\end{example}

\begin{solution}[5.2cm]
Here
\[
 A^{-1}=\begin{pmatrix}-0.1&0.3\\0.4&-0.2\end{pmatrix}.
\]
The first scaled iterate is approximately
\[
 \vect{x}_1=(0.7708,-1)\T,
 \qquad \lambda^{(1)}\approx-1.9952.
\]
The second is
\[
 \vect{x}_2=(-0.7419,1)\T,
 \qquad \lambda^{(2)}\approx-2.0020.
\]
The estimated percentage change is about $0.34\%$.  Thus the smaller eigenvalue is approximately $-2.0020$.
\end{solution}

To find an eigenvalue near a prescribed shift $a$, apply inverse iteration to
$A-aI$.  If $\beta$ is the dominant eigenvalue of $(A-aI)^{-1}$, then
\[
 \lambda\approx a+\frac1\beta.
\]

\begin{example}
Use two shifted inverse power iterations with $a=6$ and
$\vect{x}_0=(1,0)\T$ to find the eigenvalue of
\[
 A=\begin{pmatrix}7&-2\\3&-1\end{pmatrix}
\]
nearest to $6$.
\end{example}

\begin{solution}[4.2cm]
Since
\[
 (A-6I)^{-1}=\begin{pmatrix}7&-2\\3&-1\end{pmatrix},
\]
the scaled iterates are approximately
\[
 \vect{x}_1=(1,0.4286)\T,
 \qquad
 \vect{x}_2=(1,0.4186)\T.
\]
At the second step the Rayleigh estimate for $(A-6I)^{-1}$ is
$\beta\approx6.16337$.  Hence
\[
 \lambda\approx6+\frac1{6.16337}=6.1622.
\]
\end{solution}

\section{Householder Reduction and the QR Algorithm}

\begin{definition}[Householder transformation]
Let $\vect w\in\R^n$ satisfy $\vect w\T\vect w=1$.  The matrix
\[
 P=I-2\vect w\vect w\T
\]
is a Householder transformation.
\end{definition}

A Householder matrix is symmetric and orthogonal:
\[
 P\T=P,
 \qquad
 P^{-1}=P.
\]
It reflects vectors across the hyperplane orthogonal to $\vect w$.  Carefully chosen Householder transformations can introduce zeros below the first subdiagonal of a symmetric matrix while preserving symmetry and eigenvalues.  Repeating the process reduces a symmetric matrix to tridiagonal form.

\begin{example}
Reduce
\[
 A=\begin{pmatrix}
4&1&-2&2\\
1&2&0&1\\
-2&0&3&-2\\
2&1&-2&-1
\end{pmatrix}
\]
to tridiagonal form using Householder transformations.
\end{example}

\begin{solution}[4.5cm]
A sequence of orthogonal similarity transformations produces
\[
 T=\begin{pmatrix}
4&-3&0&0\\[1mm]
-3&\frac{10}{3}&-\frac53&0\\[1mm]
0&-\frac53&-\frac{33}{25}&\frac{68}{75}\\[1mm]
0&0&\frac{68}{75}&\frac{149}{75}
\end{pmatrix}.
\]
Because every step has the form $PAP\T$, the tridiagonal matrix $T$ is similar to $A$ and has the same eigenvalues.
\end{solution}

For a symmetric tridiagonal matrix, the QR algorithm repeatedly performs
\[
 A^{(k)}=Q^{(k)}R^{(k)},
 \qquad
 A^{(k+1)}=R^{(k)}Q^{(k)}.
\]
Since
\[
 A^{(k+1)}=(Q^{(k)})\T A^{(k)}Q^{(k)},
\]
all iterates are similar and have the same eigenvalues.  Under suitable conditions the off-diagonal entries tend to zero, and the diagonal entries approach the eigenvalues.

\begin{example}
Perform one QR iteration for
\[
 A=\begin{pmatrix}
3&1&0\\
1&3&1\\
0&1&3
\end{pmatrix}.
\]
\end{example}

\begin{solution}[6.0cm]
One QR factorisation is $A=QR$, with
\[
 Q\approx\begin{pmatrix}
0.94868&-0.29409&0.11625\\
0.31623&0.88226&-0.34874\\
0&0.36761&0.92998
\end{pmatrix},
\]
\[
 R\approx\begin{pmatrix}
3.16228&1.89737&0.31623\\
0&2.72029&1.98560\\
0&0&2.43270
\end{pmatrix}.
\]
Therefore
\[
 A^{(2)}=RQ\approx
\begin{pmatrix}
3.60000&0.86023&0\\
0.86023&3.12973&0.89753\\
0&0.89753&2.27027
\end{pmatrix}.
\]
This matrix is symmetric and similar to $A$, but its off-diagonal entries have begun to change toward a more diagonal form.
\end{solution}

\phantomsection
\section*{Tutorial 2}
\addcontentsline{toc}{section}{Tutorial 2}
